\section{Технологическая часть}
\label{sec:tech}

В этой части описаны три аспекта реализации: руководство пользователя (что поставить,
как собрать проект под Apache~Maven и как запустить интерпретатор), устройство
центральной структуры данных~--- раскрученной двусторонней очереди на запечатанных
типах Java~--- и цикл нормализации, в котором поле зрения обрабатывается на двух
стеках. Руководство идёт первым, чтобы запуск системы можно было воспроизвести ещё
до разбора внутренностей; язык реализации~--- Java~25~LTS~\cite{java_se25_spec}.

\subsection{Инструкция пользователя}
\label{subsec:user-manual}

Чтобы собрать и запустить интерпретатор, на машине должны стоять JDK, Maven и (только для
построения графиков) Python~3; конкретные версии перечислены в
таблице~\ref{tab:requirements}. Установленные версии удобно проверить командами
\prog{java -\/\/-version} и \prog{mvn -\/\/-version}.

\begin{table}[!htb]
\caption{\centering Требования к среде разработки}
\label{tab:requirements}
\centering
\small
\begin{tabular}{|l|c|l|}
\hline
\textbf{Компонент} & \textbf{Минимальная версия} & \textbf{Примечание} \\
\hline
Java Development Kit & 25 LTS & Сборка и запуск \\
\hline
Apache Maven & 3.9.x & Сборка проекта \\
\hline
Python~3 & 3.9+ & Только для построения графиков \\
\hline
\end{tabular}
\end{table}

Проект собирается командой Apache~Maven из корня репозитория \prog{refal5q-impl/}:

\begin{lstlisting}[language=bash,caption={Команда сборки},label={lst:build}]
mvn package -DskipTests
\end{lstlisting}

Результат сборки --- файл \prog{target/refal5q.jar}, самодостаточный
JAR-архив с прописанным классом \prog{Main-Class: refal.Main}.
Все зависимости рантайма включены в него через плагин Maven~Shade,
поэтому для запуска достаточно стандартной JVM без дополнительных библиотек.

Сборка с полным прогоном тестов выполняется командой \prog{mvn package}
(без флага \prog{-DskipTests}) или \prog{mvn verify}, которая дополнительно
формирует отчёт о покрытии кода инструментом JaCoCo.

Собранный интерпретатор запускается стандартной командой \prog{java -jar}:

\begin{lstlisting}[language=bash,caption={Форма запуска},label={lst:run}]
java -jar target/refal5q.jar [OPTIONS] program.ref
\end{lstlisting}

Позиционный аргумент --- путь к исходному файлу \prog{.ref}.
Поддерживаемые опции перечислены в таблице~\ref{tab:options}.

\begin{table}[!htb]
\caption{\centering Опции командной строки}
\label{tab:options}
\centering
\small
\begin{tabular}{|l|p{8.5cm}|}
\hline
\textbf{Опция} & \textbf{Описание} \\
\hline
\prog{--entry <имя>} &
    Явное указание входной функции; имеет приоритет над \prog{\$ENTRY} в файле \\
\hline
\prog{--step-limit <N>} &
    Максимальное число шагов нормализации (по умолчанию~--- 100\,000) \\
\hline
\prog{--expr <терм...>} &
    Начальное выражение, передаваемое входной функции вместо пустого \\
\hline
\end{tabular}
\end{table}

Код завершения процесса отражает результат выполнения
(таблица~\ref{tab:exit-codes}).

\begin{table}[!htb]
\caption{\centering Коды завершения процесса}
\label{tab:exit-codes}
\centering
\small
\begin{tabular}{|c|p{10cm}|}
\hline
\textbf{Код} & \textbf{Условие} \\
\hline
0 & Успешное завершение или \prog{<Exit 0>} \\
\hline
1 & Ошибка лексического, синтаксического или семантического анализа \\
\hline
2 & Ошибка времени выполнения: нет подходящего правила, превышение лимита шагов \\
\hline
$n$ & Значение $n$, переданное встроенной функции \prog{<Exit~n>} \\
\hline
\end{tabular}
\end{table}

Пример запуска программы \prog{hello.ref} с явным указанием точки входа:

\begin{lstlisting}[language=bash,caption={Пример запуска},label={lst:run-example}]
java -jar target/refal5q.jar --entry Go hello.ref
\end{lstlisting}

Полный набор тестов запускается одной командой Maven:

\begin{lstlisting}[language=bash,caption={Запуск тестов},label={lst:test}]
mvn test
mvn verify
mvn test -Dtest=AutotestsRunnerTest
\end{lstlisting}

Первые две команды прогоняют 143~JUnit-теста, покрывающих все компоненты
проекта. Третья команда запускает отдельный класс \prog{AutotestsRunnerTest},
который вызывает интерпретатор на 81~внешней программе из двух наборов:
автотесты репозитория \prog{refal-5j} и десахаренные тесты репозитория
\prog{refal-5-framework}. Каждая программа выполняется в отдельном процессе
с тайм-аутом 20~секунд.

Отчёт JaCoCo формируется в каталоге \prog{target/site/jacoco/}; открыть
его можно браузером, указав путь к \prog{index.html}.

%%% ─────────────────────────────────────────────────────────
\subsection{Структура данных: раскрученная двусторонняя очередь}
\label{subsec:deque-impl}

Объектное выражение в нашей реализации хранится в раскрученной (bootstrapped)
двусторонней очереди по Окасаки~\cite{okasaki_pfds}. У неё амортизированно
константная конкатенация~--- та самая операция, на которой классические
списочные интерпретаторы Рефала (включая Refal-5~PZ и Refal-5J) тратят
линейное время.

Контейнер выражения объявлен как обобщённый запечатанный интерфейс \prog{Expr<T>}
с двумя реализациями:

\begin{lstlisting}[language=Java,caption={Запечатанный интерфейс Expr<T>},label={lst:expr-iface}]
public sealed interface Expr<T> extends Iterable<T>
        permits ShallowDeque, DeepDeque {
    boolean isEmpty();
    int size();
    PeekResult<T> separateFront();
    PeekResult<T> separateBack();
    Expr<T> pushFront(T x);
    Expr<T> pushBack(T x);
    static <T> Expr<T> empty()        { return ShallowDeque.empty(); }
    static <T> Expr<T> singleton(T x) { return ShallowDeque.<T>empty().pushBack(x); }
    static <T> Expr<T> concat(Expr<T> a, Expr<T> b) { return ExprOps.concat(a, b); }
}
\end{lstlisting}

Слово \prog{sealed} запрещает любые другие наследники. Компилятор
Java~25 на этом основании проверяет полноту разбора в
\prog{switch}-выражениях, поэтому в коде конкатенации мы можем
перечислить ровно четыре случая по типам операндов и не писать ветку
\prog{default}: её отсутствие гарантировано системой типов.

Поведение обеих реализаций задаётся двумя константами.
Ёмкость мелкого буфера~--- \prog{ShallowDeque.CAPACITY = 9},
нижняя граница заполненности \prog{MIN\_FILL = 5}, то есть
$\lceil CAPACITY / 2 \rceil$. Нечётное значение выбрано по
арифметическим соображениям: при переполнении мелкого буфера суммарно
получается $CAPACITY + 1 = 10$ элементов; они делятся ровно пополам
на два фрагмента по 5, и оба фрагмента сразу удовлетворяют
\prog{MIN\_FILL}. С чётной ёмкостью пришлось бы либо мириться с
остатком в один элемент, либо добавлять отдельную обработку этого
случая~--- ни то ни другое не даёт выигрыша по свойствам структуры.

На практике глубина рекурсии \prog{DeepDeque} при таком \prog{CAPACITY}
почти не выходит за два уровня: в тестовых программах с выражениями
длиной до десятков тысяч термов средняя очередь сама ни разу не успевает
стать глубокой.

Мелкий уровень~--- это \prog{ShallowDeque<T>}, иммутабельный класс с тремя полями:
массив фиксированной длины \prog{Object[] buf} размером \prog{CAPACITY},
индекс головы \prog{head} и текущий размер \prog{sz}.
Доступ по логическому индексу~$i$ выполняется через единственную формулу:

\begin{lstlisting}[language=Java,caption={Циклическое индексирование в ShallowDeque},label={lst:shallow-idx}]
int idx(int i) {
    return Math.floorMod(head + i, CAPACITY);
}
\end{lstlisting}

\prog{floorMod} здесь стоит вместо обычного \prog{\%}, потому что в Java
оператор \prog{\%} при отрицательном делимом возвращает отрицательный
остаток, а нам нужен индекс в массиве. \prog{pushFrontTight} как раз
сдвигает голову на~$-1$, и без \prog{floorMod} в этой точке пришлось бы
вручную добавлять \prog{CAPACITY}.

В классе определены два варианта операций добавления:
\begin{itemize}
  \item \prog{pushFrontTight(T)} и \prog{pushBackTight(T)}~--- внутренние, возвращают
        \prog{ShallowDeque<T>} и бросают \prog{IllegalStateException} при переполнении;
        применяются там, где вызывающий код гарантирует наличие свободного места
        (например, при формировании буферов во время конкатенации);
  \item \prog{pushFront(T)} и \prog{pushBack(T)}~--- публичные (методы
        интерфейса \prog{Expr<T>}); при переполнении автоматически переходят
        в глубокий уровень через \prog{DeepDeque.mkDeep}.
\end{itemize}

Глубокий уровень~--- \prog{DeepDeque<T>}; он повторяет в коде ту самую тройку
$\langle front, mid, back \rangle$, что описана в п.~\ref{sec:bootstrapped-deque}:

\begin{lstlisting}[language=Java,caption={Поля DeepDeque<T>},label={lst:deep-fields}]
public final class DeepDeque<T> implements Expr<T> {
    private final ShallowDeque<T>          front;
    private final Expr<ShallowDeque<T>>    mid;
    private final ShallowDeque<T>          back;
    private final int                      sz;
    ...
}
\end{lstlisting}

Вся «раскрутка» спрятана в одном типе: поле \prog{mid} объявлено как
\prog{Expr<ShallowDeque<T>>}, а не как конкретный класс. Поэтому средняя очередь
хранит мелкие буферы, но сама подчиняется тому же интерфейсу и может на следующем
уровне снова оказаться \prog{DeepDeque}. Именно эта рекурсия по типу и даёт два
(на практике~--- редко больше) уровня вложенности без какого-либо особого кода для
каждого из них.

Конструктор \prog{mkDeep} автоматически схлопывает глубокую очередь обратно
в мелкую, когда суммарное число элементов не превышает \prog{CAPACITY}:

\begin{lstlisting}[language=Java,caption={Конструктор mkDeep с откатом в ShallowDeque},label={lst:mkdeep}]
static <T> Expr<T> mkDeep(ShallowDeque<T> front,
                           Expr<ShallowDeque<T>> mid,
                           ShallowDeque<T> back, int total) {
    if (total <= CAPACITY) {
        ShallowDeque<T> flat = ShallowDeque.empty();
        for (T x : iterableOf(front))  flat = flat.pushBackTight(x);
        for (ShallowDeque<T> bucket : iterableOf(mid))
            for (T x : iterableOf(bucket)) flat = flat.pushBackTight(x);
        for (T x : iterableOf(back))   flat = flat.pushBackTight(x);
        return flat;
    }
    return new DeepDeque<>(front, mid, back, total);
}
\end{lstlisting}

Параметр \prog{total} передаётся явно, чтобы не пересчитывать его внутри:
\prog{mid.size()} возвращает число буферов в средней очереди, а не суммарное
число листовых элементов, поэтому наивная формула
\prog{front.size() + mid.size() + back.size()} оказалась бы некорректной.
Каждый вызов \prog{mkDeep} обновляет \prog{total} инкрементально~---
прибавляя или вычитая единицу по ходу операции добавления или извлечения.

Главная операция, ради которой выбрана эта структура,~--- конкатенация. Она
реализована в служебном классе \prog{ExprOps} и разбирает четыре случая по
мелкости/глубине операндов:

\begin{lstlisting}[language=Java,caption={Разбор четырёх случаев конкатенации},label={lst:concat-cases}]
static <T> Expr<T> concat(Expr<T> a, Expr<T> b) {
    if (a.isEmpty()) return b;
    if (b.isEmpty()) return a;
    if (a instanceof ShallowDeque<T> sa && b instanceof ShallowDeque<T> sb)
        return concatSS(sa, sb);
    if (a instanceof ShallowDeque<T> sa && b instanceof DeepDeque<T> db)
        return concatSD(sa, db);
    if (a instanceof DeepDeque<T> da && b instanceof ShallowDeque<T> sb)
        return concatDS(da, sb);
    DeepDeque<T> da = (DeepDeque<T>) a;
    DeepDeque<T> db = (DeepDeque<T>) b;
    return concatDD(da, db);
}
\end{lstlisting}

В случае \emph{Deep\,+\,Deep} пограничные буферы \prog{a.back} и
\prog{b.front} склеиваются в один буфер, а \prog{concat} вызывается
рекурсивно уже на средних очередях. На каждом уровне рекурсии работа
ограничена константой, а сама рекурсия идёт не глубже логарифма от
длины выражения~--- так Окасаки и получает амортизированную
константу~\cite{okasaki_pfds}.

%%% ─────────────────────────────────────────────────────────
\subsection{Цикл нормализации}
\label{subsec:runtime-impl}

\subsubsection{Поле зрения и иерархия ViewTerm}

Нормализация выражения выполняется итеративным проходом по
\emph{полю зрения}~--- линейной последовательности элементарных единиц.
Каждая такая единица описывается запечатанным интерфейсом \prog{ViewTerm}:

\begin{lstlisting}[language=Java,caption={Иерархия ViewTerm},label={lst:viewterm}]
public sealed interface ViewTerm
        permits ViewTerm.Passive,  ViewTerm.BulkPassive,
                ViewTerm.AngOpen,  ViewTerm.AngClose,
                ViewTerm.ParenOpen, ViewTerm.ParenClose {

    record Passive(Term term)          implements ViewTerm {}
    record BulkPassive(Expr<Term> items) implements ViewTerm {}
    record AngOpen(String fn)          implements ViewTerm {}
    record AngClose()                  implements ViewTerm {}
    record ParenOpen()                 implements ViewTerm {}
    record ParenClose()                implements ViewTerm {}
}
\end{lstlisting}

Термы-атомы (\prog{Sym}, \prog{Num}, \prog{Str}) и скобочные термы \prog{Br}
после нормализации вложенных выражений представляются как \prog{Passive}.
Вызов \prog{<f args>} разворачивается в последовательность
\prog{AngOpen(f)}, элементы аргументов, \prog{AngClose}.
Скобочный терм~\prog{(e.X)} аналогично разворачивается через
\prog{ParenOpen} и \prog{ParenClose}.
\prog{BulkPassive} отвечает за оптимизацию связывания переменных и описан
в следующем подразделе.

\subsubsection{Двустековая реализация нормализатора}

Нормализация работает с двумя стеками~--- \prog{left} и \prog{right}:
\begin{itemize}
  \item \prog{right} содержит ещё не обработанные элементы поля зрения;
  \item \prog{left} накапливает уже пассивные элементы~--- те, в которых
        не осталось ни одного активного вызова.
\end{itemize}

Исходное выражение сначала сплющивается в \prog{right} слева-направо.
Затем главный цикл извлекает элементы из \prog{right} по одному с головы.
Пассивные элементы (\prog{Passive}, \prog{BulkPassive}) и открывающие
маркеры (\prog{AngOpen}, \prog{ParenOpen}) просто перекладываются в
\prog{left}. Обработка наступает при встрече закрывающего маркера:

\begin{lstlisting}[language=Java,caption={Обработка закрывающего маркера вызова функции},label={lst:angclose}]
case ViewTerm.AngClose ignored -> {
    bumpStep();
    Expr<Term> argsExpr = Expr.empty();
    String fn = null;
    while (true) {
        ViewTerm top = left.pollLast();
        if (top instanceof ViewTerm.AngOpen ao) { fn = ao.fn(); break; }
        if (top instanceof ViewTerm.Passive pt)
            argsExpr = argsExpr.pushFront(pt.term());
        else if (top instanceof ViewTerm.BulkPassive bp)
            argsExpr = Expr.concat(bp.items(), argsExpr); // O(1)
    }
}
\end{lstlisting}

При встрече \prog{AngClose} мы снимаем элементы с хвоста \prog{left}
до парного \prog{AngOpen} и попутно собираем выражение аргументов
\prog{argsExpr}. Одиночный \prog{Passive} присоединяется через
\prog{pushFront} за~$O(1)$, а \prog{BulkPassive}~--- через
\prog{Expr.concat(bp.items(), argsExpr)}, тоже за~$O(1)$
амортизированно.

После сборки аргументов для встроенной функции результат сплющивается и
помещается обратно в \prog{right}; для пользовательской функции находится
первое подходящее правило и его правая часть, с подставленными связываниями
переменных, укладывается на \prog{right} методом
\prog{pushResultToStack}.

Закрывающая скобка \prog{ParenClose} обрабатывается по той же схеме сборки с
хвоста \prog{left}, но проще: имя функции искать не нужно, а собранные термы не
уходят в вычисление, а сразу заворачиваются в скобочный терм \prog{Br}:

\begin{lstlisting}[language=Java,caption={Сборка скобочного терма по ParenClose},label={lst:parenclose}]
case ViewTerm.ParenClose ignored -> {
    Expr<Term> brItems = Expr.empty();
    while (true) {
        ViewTerm top = left.pollLast();
        if (top == null)
            throw new RefalException("Unmatched ')' in view field");
        if (top instanceof ViewTerm.ParenOpen) break;
        if (top instanceof ViewTerm.Passive pt)
            brItems = brItems.pushFront(pt.term());
        else if (top instanceof ViewTerm.BulkPassive bp)
            brItems = Expr.concat(bp.items(), brItems);
    }
    left.addLast(new ViewTerm.Passive(new Br(brItems)));
}
\end{lstlisting}

Здесь видна симметрия двух закрывающих маркеров. \prog{AngClose} и
\prog{ParenClose} оба сматывают \prog{left} до своего парного открывающего
маркера и оба используют один и тот же приём накопления (\prog{pushFront} для
одиночного терма, \prog{Expr.concat} для \prog{BulkPassive}). Различие
только в том, что получается на выходе: \prog{AngClose} отдаёт собранные
термы как аргументы вызова и запускает редукцию, а \prog{ParenClose} оборачивает
их в один \prog{Passive} с термом \prog{Br} и возвращает обратно на \prog{left}.
Готовый скобочный терм с этого момента пассивен и больше в нормализацию не
вовлекается~--- ровно поэтому вложенные скобки нормализуются изнутри наружу:
внутренний \prog{ParenClose} успевает свернуться раньше, чем до своего хвоста
доходит внешний.

Отдельно стоит отметить защиту от рассогласованного поля зрения. Если при сборке
скобок стек \prog{left} опустеет раньше, чем встретится \prog{ParenOpen}, это
значит, что во входе была лишняя закрывающая скобка~--- и нормализатор бросает
\prog{RefalException} с понятным сообщением вместо того, чтобы молча выдать
неверный результат. Симметричная проверка стоит и в \prog{AngClose}. На корректно
разобранной программе такие ситуации недостижимы (баланс скобок и угловых скобок
гарантируется ещё синтаксическим анализом), но в коде нормализатора они оставлены
как явные инварианты: при будущих изменениях нарушение баланса проявится сразу и
в точке возникновения, а не где-то ниже по вычислению.

Цикл устроен итеративно, а не рекурсивно, и это снимает зависимость от размера
стека JVM: сколь угодно глубоко вложенная рефал-программа обрабатывается на
постоянной глубине Java-стека, а расход памяти определяется только длиной поля
зрения.

\subsubsection{Оптимизация BulkPassive}
\label{subsubsec:bulkpassive}

Без оптимизации подстановка значения переменной в правой части правила
требовала поэлементной укладки: каждый \prog{Term} из связанного значения
помещался на стек \prog{right} как отдельный объект \prog{Passive}:

\begin{lstlisting}[language=Java,caption={Наивная укладка переменной (до оптимизации)},label={lst:var-naive}]
case refal.ast.Var v -> {
    Expr<Term> val = bindings.get(v.key());
    (*@{\color{black}\ttfamily // Цикл O(N): каждый терм оборачивается в Passive}@*)
    for (Term t : val) {
        right.addFirst(new ViewTerm.Passive(t));
    }
}
\end{lstlisting}

При следующей встрече \prog{AngClose} те же $N$ элементов снимались
с \prog{left} по одному. Получается ровно проблема копирования из
п.~\ref{sec:flat-list}: на программе, где аккумулятор растёт от шага к шагу,
поэлементная укладка превращает линейный по смыслу шаг в квадратичную суммарную работу.

Чинится это одним изменением: вместо перебора элементов
переменной её значение кладётся на стек одним объектом \prog{BulkPassive}:

\begin{lstlisting}[language=Java,caption={Укладка переменной через BulkPassive (O(1))},label={lst:var-bulk}]
case refal.ast.Var v -> {
    Expr<Term> val = bindings.get(v.key());
    if (!val.isEmpty()) {
        right.addFirst(new ViewTerm.BulkPassive(val)); // O(1)
    }
}
\end{lstlisting}

Когда обработчик \prog{AngClose} снимает этот объект с \prog{left},
он подмешивает его в накапливаемые аргументы вызовом \prog{Expr.concat}~---
тоже амортизированно $O(1)$ \refAlgo{lst:angclose}.

Семантику это не трогает: \prog{BulkPassive} встаёт в поле зрения ровно туда, где
раньше лежали бы те же \prog{Passive} в том же порядке. Меняется не порядок
вычисления, а только способ хранения~--- одна ссылка вместо $N$ объектов.

\subsubsection{Укладка правой части правила}

Когда сопоставление выбрало правило, на правый стек нужно положить его правую
часть~--- но не как готовое значение, а так, чтобы оставшиеся в ней вызовы функций
снова попали в нормализацию. Этим занимается \prog{pushResultToStack}: он идёт по
термам правой части справа налево (порядок важен, потому что стек разворачивает
последовательность) и для каждого терма вызывает \prog{pushTermToStack}.

\begin{lstlisting}[language=Java,caption={Укладка терма правой части на стек (ветви Var и Br)},label={lst:push-term}]
switch (t) {
    case Var v -> {
        Expr<Term> val = bindings.get(v.key());
        if (!val.isEmpty())
            right.addFirst(new ViewTerm.BulkPassive(val));
    }
    case Br br -> {
        right.addFirst(ViewTerm.ParenClose.INSTANCE);
        List<Term> items = br.items().toList();
        for (int j = items.size() - 1; j >= 0; j--)
            pushTermToStack(items.get(j), bindings, right);
        right.addFirst(ViewTerm.ParenOpen.INSTANCE);
    }
    default -> right.addFirst(new ViewTerm.Passive(t));
}
\end{lstlisting}

Каждый вид терма раскрывается по-своему. Переменная заменяется своим связанным
значением: непустое значение кладётся на стек одним объектом \prog{BulkPassive}.
Скобочный терм \prog{(\ldots)} оборачивается маркерами
\prog{ParenOpen}/\prog{ParenClose}, между которыми его содержимое рекурсивно
укладывается тем же методом справа налево. Ветвь для вызова \prog{<f \ldots>} в
листинге опущена~--- она устроена симметрично \prog{Br}: вместо скобочных маркеров
ставятся \prog{AngOpen(f)} и \prog{AngClose}, между ними рекурсивно укладываются
аргументы. Именно поэтому вложенные вызовы из правой части снова доходят до своего
\prog{AngClose} и нормализуются. Атом попадает в ветвь \prog{default} и кладётся
одним \prog{Passive}. Так одна и та же укладка обслуживает и пассивные данные, и
новые активные вызовы, не разделяя их на два разных пути.

Отдельного пояснения требует порядок обхода. И \prog{pushResultToStack}, и обе
рекурсивные ветви идут по термам \emph{справа налево} (циклы со счётчиком от
\prog{size()-1} к нулю), а каждый терм кладётся на голову правого стека через
\prog{addFirst}. Стек по своей природе разворачивает последовательность, поэтому
двойной разворот~--- обход с конца плюс вставка в начало~--- в итоге сохраняет
исходный порядок термов в поле зрения. Если бы укладка шла слева направо, правая
часть правила попала бы на стек в обратном порядке, и результат вычисления
оказался бы зеркальным. Это тонкое место: для линейного представления порядок
вставки очевиден, а на двустековой модели его приходится держать в голове, иначе
ошибка проявится не падением, а молча искажённым результатом~--- что заметно
труднее отлаживать.

Чтобы зацикленная или слишком долгая программа не висла бесконечно, каждый вход в
\prog{AngClose} проходит через счётчик шагов:

\needspace{9\baselineskip}
\begin{lstlisting}[language=Java,caption={Контроль числа шагов нормализации},label={lst:bumpstep}]
private void bumpStep() {
    steps++;
    if (steps > stepLimit)
        throw new RefalException("Step limit exceeded (" + stepLimit + ")");
}
\end{lstlisting}

Один шаг~--- это одна редукция вызова функции. Входной вызов \prog{Go} в счёт не
идёт (счётчик стартует с $-1$), поэтому \prog{<Step>} в программе и число шагов в
сообщении об ошибке совпадают с интуитивным «сколько раз сработало правило».
Порог \prog{stepLimit} задаётся опцией \prog{--step-limit}; по умолчанию~---
100\,000, чего хватает тестовым программам и заведомо мало для случайного
зацикливания.
